% Draft of report sections 4 and 5 (Olivia, 16 Nov). Paste into Overleaf.
% Needs \usepackage{graphicx}. Upload the two charts from outputs/6_results/.
% Every number comes from outputs/6_results/ (python -m keyboard_tax.analyze).
% Challenge-set numbers: review file of 16 Nov (187 items). TODO: Thomas confirms that he
% checked the rows himself (see notes/olivia_steps_4-6_notes.md), otherwise soften "checked by our native speaker".

\section{Translation systems}

We train no model. Every version of every sentence is translated by three freely available
systems that represent three kinds of translator:
\textbf{NLLB-200} (distilled, 600M parameters), a multilingual translation model for 200 languages;
\textbf{Opus-MT ROMANCE$\rightarrow$en}, a small Marian model trained on OPUS data for a group of
Romance languages (there is no official Portuguese--English Opus-MT model); and
\textbf{Qwen2.5-7B-Instruct}, an open-weights instruction-tuned LLM. We did not use a Claude model,
because Claude suggested the first-pass labels of our word list.

NLLB and Opus-MT use beam search (beam size 4, at most 256 new tokens). The LLM gets one sentence
per request with one fixed prompt (``Translate the following Portuguese sentence into English.
Output only the English translation, with no notes or explanations.''), no sampling, and 4-bit
weights; fixed rules remove extra text such as ``Here is the translation:''. The prompt was not tuned
on \texttt{devtest}. All models ran on a free Google Colab GPU. Each distinct text is translated once
per system: 3,290 versions $\times$ 3 systems give 9,870 translations, none of them empty. For 4
sentences the LLM answered partly in Chinese; we keep these answers and show in Section~5 (Robustness) that they do not change the results.

\section{Evaluation}

\paragraph{Scores.} We compute BLEU and chrF with sacreBLEU\footnote{BLEU
\texttt{nrefs:1|case:mixed|eff:no|tok:13a|smooth:exp|version:2.6.0}; chrF
\texttt{nrefs:1|case:mixed|eff:yes|nc:6|nw:0|space:no|version:2.6.0}.} and COMET
(\texttt{wmt22-comet-da}), our main score because it reacts to meaning. COMET always gets the
\emph{clean} Portuguese sentence as its source, so it judges each translation against what the writer
meant. For every sentence we take the \emph{drop}: score of the clean translation minus score of the
damaged one (positive = quality lost). Comparisons are paired, sentence by sentence, and we report
the mean drop with a 95\% confidence interval (CI). Class~A only, Class~B only and typos exist for 422
sentences, so they are compared on those 422 sentences only. Both versions strip 1.27 words per
sentence on average, and the typo version changes as many letters as the Class~B version.

\paragraph{Targeted check.} Scores average over whole sentences, so we also check single words. Our
challenge set has 187 items (170 sentences, 74 Class~B words) in which a lost accent changes the English
meaning, e.g.\ \emph{país} ``country'' $\rightarrow$ \emph{pais} ``parents''. Each item lists English
words for the right and for the wrong reading. A translation counts as \emph{right} if it contains a
right-reading word and no wrong-reading word, \emph{wrong} in the opposite case, and \emph{unclear}
otherwise. Only items where the word really lost its accent are counted. We drafted 159 items from a hand-written list of word pairs; the review by our native speaker kept
136, corrected 13, removed 10 and added 38. Very common pairs such as \emph{é/e} (``is''/``and'') are
included only where the English reference allows a word match; the COMET drop covers the rest.

\begin{table}[t]
\centering\small
\setlength{\tabcolsep}{4pt}
\begin{tabular}{lccc}
\hline
 & NLLB & Opus-MT & LLM \\
\hline
COMET, clean (1,012 sentences) & 0.890 & 0.882 & 0.891 \\
\hline
\multicolumn{4}{l}{\emph{COMET drop $\times$100, mean [95\% CI]}} \\
All accents removed (1,012) & 1.10 [0.90, 1.30] & 2.24 [1.95, 2.53] & 0.30 [0.17, 0.43] \\
Class A only (422) & 0.25 [0.09, 0.40] & 0.35 [0.17, 0.52] & 0.05 [$-$0.03, 0.14] \\
Class B only (422) & 0.70 [0.45, 0.96] & 1.73 [1.36, 2.09] & 0.26 [0.13, 0.40] \\
Typos, same size (422) & 1.81 [1.45, 2.17] & 4.37 [3.84, 4.91] & 0.70 [0.45, 0.95] \\
B minus A (422) & 0.46 [0.19, 0.73] & 1.38 [0.98, 1.78] & 0.21 [0.06, 0.36] \\
B minus A, without \emph{é} (240) & 0.17 [$-$0.14, 0.48] & 0.83 [0.34, 1.32] & $-$0.12 [$-$0.25, 0.00] \\
\hline
BLEU, clean $\rightarrow$ all removed & 48.2 $\rightarrow$ 46.5 & 45.2 $\rightarrow$ 41.8 & 45.8 $\rightarrow$ 44.9 \\
Challenge set, share right, same & 95\% $\rightarrow$ 83\% & 99\% $\rightarrow$ 75\% & 97\% $\rightarrow$ 91\% \\
\hline
\end{tabular}
\caption{Main results. A drop is the clean score minus the damaged score of the same sentence.}
\label{tab:results}
\end{table}

\begin{figure}[t]
\centering
\includegraphics[width=0.85\linewidth]{chart_drops.png}
\caption{COMET drop per kind of damage on the 422 sentences with all five versions (95\% CI).}
\label{fig:drops}
\end{figure}

\paragraph{RQ1: How much do missing accents cost?} Removing all accents lowers COMET for every system
(H1 supported; Table~\ref{tab:results}). Opus-MT loses most, twice as much as NLLB; the LLM loses
about a quarter of what NLLB loses. Corpus BLEU falls by 0.9 to 3.4 points.

\paragraph{RQ2: Is Class B worse than Class A, and than typos?} With the same number of words stripped,
Class~B costs more than Class~A for all three systems (H2 supported; the CI of ``B minus A'' is above
0). Without the sentences whose stripped Class~B word is \emph{é}, the difference holds only for
Opus-MT (H2b): for NLLB and the LLM, most of the Class~B cost comes from \emph{é} turning into
\emph{e} (``and''). Typos of the same size cost more than Class~B, for every system, by a factor of
about 2.5 (H3 not supported, Figure~\ref{fig:drops}). A likely reason: a stripped accent mostly leaves
a real word or a common spelling the models have seen, while a typo leaves a form they have not.

As planned, we also fitted one regression per system,
$\Delta = \beta_0 + \beta_1 b + \beta_2 a + \beta_3 c$, on all damaged versions with errors clustered by
sentence ($b$, $a$: Class~B and Class~A words stripped; $c$: letters changed). It does not work on our
data. A stripped Class~A word changes almost always one letter, so $a$ and $c$ move together
(correlation 0.88), and the fit cannot tell them apart: $\beta_2$ comes out \emph{negative}
($-$0.018 for Opus-MT) and $\beta_3$ positive, which only shifts the same cost between two near-identical
predictors. H4 ($\beta_1 > \beta_2$, $\beta_3 \approx 0$) therefore cannot be decided; the paired
comparisons above are our main test.

\begin{figure}[t]
\centering
\includegraphics[width=0.85\linewidth]{chart_targeted.png}
\caption{Share of challenge items translated with the right reading (only items whose word lost its accent).}
\label{fig:targeted}
\end{figure}

\paragraph{RQ3: Do the scores notice?} On the 170 challenge sentences, the share of right readings falls
by 12\% [7, 17] for NLLB, 24\% [18, 30] for Opus-MT and 6\% [2, 10] for the LLM when all accents are
removed (Figure~\ref{fig:targeted}). BLEU on the same sentences falls by only 5\%, 8\% and 3\%. For NLLB
and Opus-MT the meaning loss is clearly larger than the BLEU loss (H5 supported); for the LLM the
difference is not certain. The kind of loss matters. Pairs that only change the word class
(\emph{análise} ``analysis'' $\rightarrow$ \emph{analise} ``analyzes'', 93 items) almost never fool a
system (92\% right or more): the context decides. Pairs that change the meaning (55 items) do: Opus-MT
keeps the right reading in 51\% of them, NLLB in 76\%, the LLM in 95\%. Tense pairs, mostly future verbs
(\emph{acontecerá} $\rightarrow$ \emph{acontecera}, 15 items), fail for every system (7--33\% right):
most translations follow the form without the accent (``had taken place'', ``would help'').

\paragraph{RQ4: Which system copes best?} The LLM. Its drop after removing all accents is 0.80 [0.57,
1.02] COMET points ($\times$100) smaller than NLLB's and 1.94 [1.64, 2.24] smaller than Opus-MT's, and
it keeps the most right readings (H6 supported).

\paragraph{Examples.} Three Class~B errors (clean source $\rightarrow$ damaged source: damaged translation):
\begin{itemize}\itemsep0pt
\item Opus-MT: ``causando as \emph{marés} [\ldots] uma \emph{força}'' $\rightarrow$ ``\emph{mares}
[\ldots] \emph{forca}'': ``causing the \emph{seas}, the Milky Way exerts a \emph{gallows} over the
galaxy'' (COMET $-$0.08).
\item NLLB: ``os animais que a jaguatirica \emph{caça}'' (hunts) $\rightarrow$ ``\emph{caca}'': ``the
animals that the jaguars \emph{poop}'' (COMET $-$0.05).
\item LLM: ``rapidamente \emph{desaparecerá}'' (will disappear) $\rightarrow$ ``\emph{desaparecera}'':
``quickly \emph{disappeared}'' (COMET $-$0.02).
\end{itemize}
The last example shows the problem of the scores: a reversed tense costs only 0.02 COMET.

\paragraph{Robustness.} Leaving out the 4 sentences with Chinese LLM answers changes
no drop by more than 0.001 and no conclusion. Leaving out the \emph{é} sentences is discussed under RQ2. Limits of this
evaluation: one native speaker checked the challenge set; a word match marks 3--10\% of the damaged
items as unclear; and FLORES+ is formal text, while people drop accents mostly in informal writing.
